\chapter{LINE Messaging API 與系統概念}

\section{LINE Messaging API 簡介}
LINE Messaging API 是 LINE 提供給開發者的訊息整合介面。開發者可以建立官方帳號的後端服務，接收使用者傳送的事件，並透過 API 回覆或主動推播訊息。它的核心價值在於把熟悉的聊天介面與自動化服務連接起來，使用者不需要另外安裝應用程式就能使用功能。

本作業實作的是 Echo Bot。使用者在 LINE 聊天室輸入文字後，LINE 平台會將事件以 HTTP POST 傳送到伺服器的 Webhook URL；Flask 解析請求並交由 SDK 驗證簽章，再以 reply token 回傳相同文字。這個例子雖然簡單，卻涵蓋正式機器人的重要基礎：認證、事件路由、回覆期限與伺服器部署。

\section{常見應用場景}
\begin{itemize}[leftmargin=2.5em]
    \item \textbf{客服與 FAQ：}將常見問題整理成關鍵字或意圖，提供即時回答。
    \item \textbf{通知與提醒：}透過 Push Message 傳送訂單狀態、課程提醒或系統告警。
    \item \textbf{活動與會員服務：}利用 Rich Menu 導向報名、查詢與會員功能。
    \item \textbf{資料查詢：}將資料庫、校務系統或其他 API 包裝成聊天操作。
\end{itemize}

\section{訊息流程與元件角色}
系統可以分成 LINE 使用者、LINE Platform、公開 Webhook 伺服器與 Messaging API 四個部分。使用者訊息先抵達 LINE Platform，平台建立事件 JSON 並送到 Webhook；伺服器驗證 \texttt{X-Line-Signature} 後，再使用 reply token 呼叫回覆 API。

\begin{figure}[H]
    \centering
    \fbox{\parbox{0.88\textwidth}{\centering
    使用者\quad $\longrightarrow$\quad LINE Platform\quad $\longrightarrow$\quad Flask Webhook\\[0.4cm]
    \hfill $\longleftarrow$ 回覆訊息 $\longleftarrow$\hfill}}
    \caption{LINE Echo Bot 的訊息往返流程}
    \label{fig:message-flow}
\end{figure}

\begin{table}[H]
    \centering
    \caption{主要元件與責任}
    \label{tab:components}
    \begin{tabular}{p{0.25\textwidth}p{0.62\textwidth}}
        \toprule
        元件 & 責任 \\
        \midrule
        LINE Platform & 管理使用者訊息、事件轉送與 Channel 認證。 \\
        Webhook URL & 提供公開 HTTPS POST 端點，接收事件資料。 \\
        Channel Secret & 用來驗證請求簽章，防止偽造事件。 \\
        Access Token & 伺服器呼叫 Messaging API 時使用的授權憑證。 \\
        Flask 應用程式 & 接收請求、交給 Handler 分派，並呼叫回覆 API。 \\
        \bottomrule
    \end{tabular}
\end{table}

\section{本文件的實作範圍}
本文件採用 Python 3、Flask、\texttt{line-bot-sdk} 第三版與 \texttt{python-dotenv}。實作範圍先聚焦在文字訊息回覆，藉此說明最重要的設定與驗證步驟；圖片訊息、Rich Menu 和資料庫可以在這個基礎上延伸。
